% RESEARCH_DRAFT_SHA256: 9d3d0037fd42d33bfe23c087be471dc5d9f7f8c897bb9f94918691c3968ba6bd
% Layout and source-reference wrappers expanded for the web renderer.
\documentclass[11pt]{article}
\usepackage[margin=25mm]{geometry}
\usepackage{amsthm}
\usepackage{fontspec,unicode-math}
\setmainfont{Latin Modern Roman}
\setmathfont{Latin Modern Math}
\usepackage{amsmath,amssymb,mathtools,xurl,hyperref}
\hypersetup{colorlinks=true,urlcolor=blue,linkcolor=blue}
\setcounter{secnumdepth}{0}
\setlength{\parindent}{0pt}
\setlength{\parskip}{5pt}
\setlength{\emergencystretch}{3em}
\newcommand{\sref}[2]{\hyperlink{source:#1:#2}{[S#2]}}
\newcommand{\cataloguescope}[1]{\paragraph{Source record.}#1}
\newcommand{\floor}[1]{\left\lfloor#1\right\rfloor}
\newtheorem{proposition}{Proposition}
\newtheorem{lemma}[proposition]{Lemma}
\begin{document}
\paragraph{Research manuscript.} OpenAI \texttt{gpt-6-astra}, reasoning \texttt{ultra}; 11 October 2026.
\paragraph{Preserved source body.} The complete catalogue mathematical text and sources follow with layout and source links expanded, including source qualifications and its historical preparation notes. New research begins at the section headed Research attempt.

\section{Decidability of Piecewise Linear Complex Embeddings}\label{statement}

\subsection{Definitions and mathematical statement}

\paragraph{Shared notation.}Standard number systems are $\mathbb N,\mathbb Z,\mathbb Q,\mathbb R,\mathbb C$. Specialized notation and index ranges follow the cited source and the definitions below.


A polyhedron is an intersection of finitely many closed affine half-spaces; a polytope is a bounded polyhedron. Its faces are its intersections with supporting hyperplanes.

A finite simplicial complex is a finite family of vertex subsets closed under taking subsets. Its dimension is one less than its largest simplex size. A piecewise-linear embedding is an injective map linear on simplices after finite subdivisions.

A continuous piecewise-linear function is continuous and agrees with an affine function on each member of a finite polyhedral partition. A difference-of-convex decomposition writes it as $g-h$ with both $g$ and $h$ convex.
\paragraph{Record scope and qualifications.}
The catalogue title is retained as metadata. Named certificates, contracts, bounds or frameworks in that title are not asserted as results of the original source extract.
Source-record limitation: Same full piecewise-linear embeddability decision target as OWR ID 30001060; that record provides a related canonical formulation.

Related catalogue records: \href{https://mehmetmars7.github.io/Erdosproblems-llm-hunter/problem.html?type=open_problems&id=30001060}{30001060}. Separate source IDs are preserved for their provenance.
\paragraph{Statement recorded in the supplied source.}
For fixed positive integers $k\leq d$ and a finite $k$-dimensional simplicial complex $K$, decide whether $K$ admits a piecewise-linear embedding into $\mathbb R^d$. \href{https://doi.org/10.4171/owr/2019/39}{[S2]}
\paragraph{Source status.}
The supplied catalogue labels this source record partially solved; that metadata was checked against the supplied snapshot on 2026-10-11. The label is preserved for provenance and is not an independent verification of present open status or mathematical correctness. The catalogue literature assessment is dated 2026-08-22.
\subsection{Short English statement}
The source asks about decidability of Piecewise Linear Complex Embeddings.
\subsection{Sources}
\begin{itemize}
\item[S1] Ulam AI and the UnsolvedMath Contributors, UnsolvedMath: A Curated Collection of Open Mathematics Problems, catalogue entry. Catalogue curation and metadata: CC BY 4.0. The source extract is retained; definitions, English summaries and qualifications are editorial additions.
\url{https://huggingface.co/datasets/ulamai/UnsolvedMath}.
\item[S2] Geometric, Algebraic, and Topological Combinatorics (2020), Oberwolfach Reports, DOI 10.4171/owr/2019/39, pp. 2395--2472
\url{https://doi.org/10.4171/owr/2019/39}.
\item[S3] Filakovsky, Wagner, and Zhechev
\url{https://doi.org/10.1137/1.9781611975994.47}.
\item[S4] Current embeddability classification
\url{https://academic.oup.com/imrn/article/2024/17/12433/7728827}.
\item Source-specific Oberwolfach Reports licences, including CC BY-SA 4.0 where applicable, are retained; see THIRD\_PARTY\_NOTICES.md.
\end{itemize}
\paragraph{Preparation and review.}
The original source and UnsolvedMath contributors retain attribution for the source extract. Codex (GPT-6; exact serving revision not exposed) assisted with definitions, English summaries, scope notes and formatting. No new human mathematical authorship or independently verified solution is asserted. Catalogue identity checked against the complete supplied registry and current repository. Distinct upstream records, including repeated questions and broad or damaged extracts, are retained for complete coverage of the source CS category. A repeated record is not a new independent conjecture. Complete mathematical-equivalence and current-literature review has not been performed. Related canonical identities: 30001060. Preparation and final review follow contribution prompts C and D, with source gaps disclosed. Structural and rendering checks do not establish a new theorem or open status.
\label{end:statement}

\section{Research attempt}
\subsection{Literature, source context and scope}
The full original catalogue body and source qualifications are preserved above. The relevant primary workshop passages were read in the cached publisher report. Historical claims and subsequent attributed results are distinguished from the mathematical work below. A restricted example, standard lemma or known theorem is not counted as new progress on the entire original target. Missing assumptions and repeated or rhetorical extracts are not silently turned into different solved problems.

\subsection{A complete general-position regime and its failed extension}
The primary report already records polynomial and undecidable regimes of piecewise-linear embeddability. We prove the elementary always-embeddable regime $d\ge2k+1$, then test whether the same placement argument extends one dimension lower. Let $K$ be a finite simplicial complex of dimension at most $k\ge1$, with vertices numbered $1,\ldots,N$.
\begin{proposition}
Mapping vertex $i$ to $(i,i^2,\ldots,i^{2k+1})$ and extending affinely over each simplex gives a piecewise-linear embedding of $K$ in $\mathbb R^{2k+1}$.
\end{proposition}
\begin{proof}
Any $r\le2k+2$ distinct image vertices are affinely independent. Indeed the first $r$ rows of their homogeneous coordinate matrix are the powers $0,1,\ldots,r-1$; its Vandermonde determinant is $\prod_{i<j}(t_j-t_i)\ne0$. This identity follows because the determinant is an alternating polynomial of degree $r(r-1)/2$, divisible by each difference, and comparison of the leading product coefficient gives constant one.
Each simplex has at most $k+1$ vertices, so its affine map is injective. The union of the vertex sets of two simplices has at most $2k+2$ elements and is affinely independent. If their image convex hulls intersect, equality of the two convex combinations and uniqueness of affine coordinates on that union force every nonshared vertex coefficient to vanish. The intersection is therefore exactly the convex hull of the shared vertices, the image of their common face. Thus the map is injective on the whole realization, and it is continuous since its affine restrictions agree on faces. The finite realization is compact, and its image lies in a Hausdorff Euclidean space; a continuous bijection from a compact space to its image has continuous inverse because compact sets have closed images. This proves embedding.
\end{proof}
For fixed $k$, each coordinate is an integer of $O(k\log(N+1))$ bits and can be produced with $O(k)$ multiplications per vertex. The same map into larger $d$ is obtained by zero-padding. This is already a known easy regime and settles no disputed pair of dimensions.
The first attempted extension uses $(i,i^2)$ for graph vertices in the plane. For $K_4$, edges joining vertices 1 to 3 and 2 to 4 lie respectively on $y=4x-3$ and $y=6x-8$. They meet at $(5/2,7)$ in the interior of both segments. Their abstract edges are disjoint, so this placement fails. This is not a nonembeddability certificate: $K_4$ embeds with three vertices at $(0,0),(2,0),(0,2)$ and the fourth at $(1/2,1/2)$, with all six edges drawn as straight segments. The three interior segments meet only at their common interior vertex and do not cross the triangle boundary.
\subsection{Verification and remaining gaps}
Finite Vandermonde determinants and the exact crossing coordinates were checked. At $d=2k$, the union of two disjoint maximal faces can have $2k+2$ vertices, more than affine independence permits in that dimension. Hence the general-position proof has an actual missing premise; changing the placement might still succeed, as the $K_4$ example shows. This does not decide the lower-dimensional unresolved regimes or reproduce the source's sophisticated obstruction and undecidability theorems.

\section{References and source attribution}
Primary workshop report: \url{https://doi.org/10.4171/owr/2019/39}. Relevant passages read: 2019/39 pp. 2458--2459, EMBED problem, known regimes and the reported undecidability theorem.

All retrievals were checked on 11 October 2026. Original question attribution remains in the preserved Sources section. External theorems are attributed and not claimed as new discoveries or fully verified merely by citation.


\paragraph{Integration review.} Codex (GPT-6; exact serving revision not exposed) reviewed the written arguments and expanded layout and source-reference wrappers. Repository metadata and the production web renderer were checked using Python 3.13 and JavaScript with MathJax 3.2.2. This model-assisted review is not independent human proof verification; no PDF compilation or formal proof checking is asserted.

\section{Completion Estimate}
\noindent\textbf{COMPLETION: 0\%}
No substantiated new progress toward the entire original target is claimed. The worked results and failed route are delimited above; a known result, illustrative model, or auxiliary obstruction is not counted as original-target completion.

\section{Final status and attribution}
\textbf{Attempt status: unresolved.} This describes this research attempt, not a new ruling on the source catalogue's status or on external literature claims.
The mathematics and this self-audit were generated by OpenAI \texttt{gpt-6-astra}, reasoning \texttt{ultra}, on 11 October 2026. Codex (GPT-6; exact serving revision not exposed) orchestrated the batch, reviewed the written mathematical arguments, and contributed formatting/repository review. This additional model review is not independent human verification. No human mathematical authorship, independent human proof checking, or formal verification is claimed. The historical preparation paragraphs in the reproduced catalogue body concern that body only.

Contribution prompt A was applied to the complete definition and supplied source record, followed by the adversarial checks of prompt B. Local reproduction files are \texttt{batch022/checks.py}, \texttt{computations.json}, and the per-ID audit JSON. They were not public links at preparation time. Numerical checks are finite and do not prove unrestricted results. No LaTeX engine was found locally; compilation was not performed by this research agent.

\end{document}
